% =====================================================================
% Methodology chapter, PRJ93.
% main.tex issues \chapter{Methods} before \input-ing this file.
%
% PREAMBLE DEPENDENCIES (main.tex does not yet carry these):
%   \usepackage[ruled,vlined,linesnumbered]{algorithm2e}
%   \usetikzlibrary{positioning,arrows.meta,fit,backgrounds,decorations.pathreplacing}
% fig:blocks (out/fig_blocks.tex) and fig:pipeline (out/fig_pipeline.tex) are
% inserted at the marked points. Neither has been compiled.
%
% Labels: every label carried by the superseded revision is preserved, so no
% inbound \ref from results.tex or conclusion.tex breaks. sec:repro, sec:mcs,
% sec:ruler-functional and sec:ruler-ellel no longer head sections; each sits on
% the passage that now carries its content.
% =====================================================================
\label{chap:methodology}

The methods below were applied to a three-venue hospitality estate in order to learn each
venue's daily trading rhythm from point-of-sale history, bound that rhythm with a calibrated
interval, and monitor departures from it. Measurement instruments are specified before the
models they grade, because on these series the choice of error denominator changes conclusions
rather than only their scale. Procedural detail (parameter grids, computational environment
and robustness variants) is in Appendices~C and~D.

\section{Study design and data sources}
\label{sec:design}

The estate comprises three trading venues in Lancashire: the Beer Hall, Ellel and Two River
Taps. Table~\ref{tab:venues} gives each frame's length and trading intensity. The three differ
enough to function as three regimes rather than three replicates, which is why several methods
below are specified per venue. Two River Taps ceased trading on 2026-05-08 and its history
cannot grow, which makes it a control: a measurement moving at the Beer Hall or Ellel while
fixed at Two River Taps is attributable to accumulated history rather than to a change in
procedure.

\begin{table}[t]
\centering
\small
% Trace: frame lengths and trading intensity from build_features after trim_to_active,
% store ceiling 2026-07-07; brain/log/45_G17d_Intermittency_and_Occurrence.md.
\caption[The three venues, their frame lengths and trading intensity at the store ceiling]{The
three venues, their frame lengths and trading intensity at the store ceiling. Frame length is
calendar days after trimming to each venue's active span, at ceiling 2026-07-07. Trading
intensity is trading days per calendar week. Two River Taps closed on 2026-05-08, so its frame
is fixed.}
\label{tab:venues}
\begin{tabular}{lrrp{0.34\textwidth}}
\toprule
Venue & Frame & Trading days/wk & Role in the design \\
\midrule
Beer Hall      & 399 & 5.30 & Anchor; richest history; methods tested here first \\
Ellel          & 386 & 1.21 & Sparse and booking-driven; constrains method choice \\
Two River Taps & 331 & 5.92 & Closed 2026-05-08; frozen control series \\
\bottomrule
\end{tabular}
\end{table}

% Trace: brain/store/manifest.json per_venue_counts; brain/config.py:110-127.
The unit of analysis is a trading venue, meaning a fixed site with an opening calendar against
which a daily rhythm can be learned and a departure from normal defined. A fourth location in
the warehouse books off-site event transactions and falls outside that boundary: it has no site
and no opening calendar, so it carries no rhythm to learn whatever volume it holds. This is a
categorical boundary on the unit of analysis rather than a threshold applied to data, and the
location's 203 records across two dates corroborate it without having set it.

The measured input is a point-of-sale export of ex-VAT daily revenue at venue, category and item
level, with all work reported at venue level unless stated. Records were trimmed to each venue's
active span and the warehouse pinned at a store ceiling of 2026-07-07, without which a
rolling-origin evaluation is not reproducible.
\label{sec:repro}
Library versions, the model revision hash and the compute device are pinned per environment and
stamped on every artefact for the same reason; Appendix~C records them.

% Trace: brain/log/51_G17j_Chatlog_Signal.md (corpus composition).
No revenue observation is discarded as an outlier. Extreme days are the object of study rather
than contamination, and removing them would remove the deviations
Section~\ref{sec:detection} exists to find. Cleaning is therefore confined to adjudicating what
counts as a trading day. Three days carry till activity and exactly zero net revenue: one is a
comped day at a venue open and serving, which is a trading day with zero takings, while a
reversed sale and a batch of voided lines are artefacts and are not trading days. A second
input, a corpus of 735 staff and assistant chat messages, supports
Section~\ref{sec:chatlog}. Stock movements and checklist completions, two further domains named
in the project specification, were unavailable; the system therefore learns from two of the four
specified domains, and the consequence is stated where it binds. Figure~\ref{fig:pipeline} traces
the path from that export
to an intervention, with the subsection number printed in each block locating the step that
governs it.

\section{Accuracy measures and the denominator basis}
\label{sec:ruler}

Accuracy must be comparable across venues trading on 5.3, 1.2 and 5.9 days a week respectively,
which rules out an unscaled error as the primary instrument. The scaled-error family of
\citet{hyndman_another_2006} divides a forecast's error by the in-sample error of a naive
benchmark. Taking the absolute form first, because it defines the denominator both measures
share, the seasonal variant used throughout scales by the in-sample mean absolute error of the
seasonal-naive method at lag $m$ \citep{hyndman_forecasting_2021},
\begin{equation}
\mathrm{MASE} \;=\; \frac{\frac{1}{h}\sum_{t=1}^{h}\lvert y_t - \hat{y}_t\rvert}{s},
\qquad
s \;=\; \frac{1}{\lvert \mathcal{D}\rvert}\sum_{t \in \mathcal{D}} \bigl\lvert y_t -
y_{t-m}\bigr\rvert,
\label{eq:mase}
\end{equation}
with $m = 7$ and $\mathcal{D}$ the in-sample index pairs entering the denominator. A value below
one beats the benchmark. The headline measure is the squared counterpart, the root mean squared
scaled error of \citet{makridakis_m5_2022}, whose scale \citet{hewamalage_look_2021} give as
\begin{equation}
s_{\mathrm{M5}} \;=\; \frac{1}{n-1}\sum_{i=2}^{n}\bigl(y_i - y_{i-1}\bigr)^2 ,
\label{eq:rmsse-m5}
\end{equation}
for $n$ the length of the observed period. It is computed on the same bases as
Equation~\ref{eq:mase} rather than on the lag-one denominator, so the two share one ruler and
differ only in loss function.

On a series with closed days $\mathcal{D}$ is not uniquely defined, and the admissible readings
disagree by enough to reverse a conclusion. Four bases were therefore named and all four
computed, with a bootstrap interval on each scale rather than a point estimate, as
Table~\ref{tab:bases} reports. At Ellel the scale runs from 180 to 806 across four admissible
readings, so the same forecast reads as 0.411 or 0.092 by that choice alone; the narrowest
reading there rests on 28 admissible pairs and carries a 95 per cent interval 65.6 per cent wide.
% Trace: brain/log/45_G17d_Intermittency_and_Occurrence.md:105-118, function bootstrap_scale,
% B = 10,000, percentile interval, as_of = 2026-07-07.

\begin{table}[t]
\centering
\small
\caption[The four denominator bases, their bootstrap scales, and the scaled error each
induces]{The four denominator bases, their bootstrap scales, and the scaled error each induces.
Percentile intervals from 10,000 resamples of the absolute lag-difference vector, as-of
2026-07-07. Pairs is the sample the denominator rests on. Dashes mark bases inadmissible at that
venue. No induced error is quoted at the dormant venue.}
\label{tab:bases}
\begin{tabular}{llrlrl}
\toprule
Venue & Basis & Scale & 95\% interval & Pairs & Induced MASE \\
\midrule
Beer Hall & \texttt{calendar\_lag7}         & 315.7 & [270.0, 365.0]  & 392 & 0.571 [0.494, 0.668] \\
Beer Hall & \texttt{calendar\_lag7\_active} & 386.9 & [332.1, 446.4]  & 276 & 0.466 [0.404, 0.543] \\
\midrule
Ellel     & \texttt{calendar\_lag7}          & 180.1 & [135.3, 229.9]   & 385 & 0.411 [0.322, 0.547] \\
Ellel     & \texttt{trading\_lag7}           & 770.8 & [604.3, 952.7]   & 61  & 0.096 [0.078, 0.122] \\
Ellel     & \texttt{trading\_same\_weekday}  & 806.2 & [639.1, 979.2]   & 67  & 0.092 [0.076, 0.116] \\
Ellel     & \texttt{calendar\_lag7\_active}  & 754.0 & [522.0, 1016.4]  & \textbf{28} & 0.098 [0.073, 0.142] \\
\midrule
Two River Taps & \texttt{calendar\_lag7\_active} & 173.2 & [148.6, 202.1] & 268 & --- \\
\bottomrule
\end{tabular}
\end{table}

\label{sec:ruler-ellel}
The ruling follows. The Beer Hall and Two River Taps are scored on
\texttt{calendar\_lag7\_active}; Ellel has no defensible scaled basis and is scored on unscaled
error, root mean squared as headline and mean absolute as secondary. That departs from the
comparability requirement this section opens with and needs an argument rather than a note. The
remedy the literature offers for demand this intermittent is
\citet{chatfield_all-zero_2007}'s, which is to stop tuning against an error measure and score
against system cost instead. That cost is the cost of an inventory system in which the forecast
\emph{is} the replenishment quantity, summing ordering, holding and shortage costs each priced
per unit. Ellel's estimand is daily revenue in pounds. Revenue is not ordered, held or
backordered, so those three rates are not unelicited here: no unit exists whose carrying or
shortage they would price, and they are undefined. The remedy is defined for a different
estimand, which is a property of the problem rather than a gap in the preparation. What does
transfer is elsewhere in the same study, whose own percentage and geometric measures had to be
modified or dropped because neither survives a demand or an error of zero, which is direct
support for dropping the denominator where the zeros are thickest. One replenishment decision does
exist
here: a stock signal flags a reorder on projected days of cover, at the Beer Hall alone, from a
product-level forecast in units per day rather than the revenue series, and carries no ordering,
holding or shortage cost.

\label{sec:ruler-functional}
Which functional a measure elicits is distinct from which scale it carries, and here it is the
sharper question. Following \citet{kolassa_we_2023}, an absolute-error measure is minimised by the
conditional median and a squared-error measure by the conditional mean, so the two reward
different forecasters rather than grading one on different scales. The estate's decision layer
acts on expected takings, so the squared measure carries the headline and the absolute measure is
retained as a labelled secondary, both because it is more widely reported and because a conclusion
holding under only one of them is a conclusion about the measure. That choice sharpens one
limitation rather than settling it: the served foundation model returns a median under a mean's
name, so the headline measure elicits a functional the served forecaster cannot produce.
Integrating over the emitted quantiles is the available remedy and was declined because it would
change the served model after the evaluation was frozen;
Section~\ref{sec:conclusion-limitations} carries the consequence. Every scaled figure in this
document states its basis and its as-of date, since the denominator of Equation~\ref{eq:mase} is
computed in-sample and grows with the warehouse. For interval quality the Winkler score is used
instead, a proper scoring rule in the units of the data and so usable where no scaled basis is
defensible.

\section{Demand-pattern classification}
\label{sec:intermittency}

Whether a series is treated as intermittent is decided on the average inter-demand interval $p$,
the mean gap between consecutive non-zero days, and the squared coefficient of variation of
demand sizes $v$, following \citet{syntetos_categorization_2005}. Two constants in that scheme
are arithmetically incorrect, and the corrections of \citet{kostenko_note_2006} are used: the
limiting interval is $p = 4/3$ rather than $1.32$, and the maximum squared coefficient of
variation $v = 0.5$ rather than $0.49$. The original pair is retained under an explicit name so
a classification can be reported under each. The same authors give a selection rule between
Croston's method and the Syntetos–Boylan approximation by dividing the classification plane
diagonally, preferring the approximation when
\begin{equation}
v \;>\; 2 - \tfrac{3}{2}p .
\label{eq:sba}
\end{equation}
Its direction is fixed by geometry: the diagonal passes through $(p,v) = (1, 0.5)$ and
$(p,v) = (4/3, 0)$, the maxima at which Croston's method can still be preferred, so Croston
occupies the region nearer the origin and the approximation is preferred above the line.

Equation~\ref{eq:sba} carries no information over the series it governs, and saying so matters
because a count of selections would otherwise read as an empirical finding. The intermittency
cutoff is $p \geq 4/3$ and $2 - \tfrac{3}{2}\cdot\tfrac{4}{3} = 0$ exactly, so the selection
threshold is non-positive for any series at or above that cutoff while $v \geq 0$ always.
Classification as intermittent therefore \emph{entails} selection of the approximation. The rule
discriminates only in the non-intermittent region, where it is not consulted, and that is a
property of the scheme rather than of this estate.

Which estimator each node adopts is a second and more consequential decision. Selecting on the
classification alone, the more literal ex-ante reading, was rejected on evidence: the intermittent
estimator loses to the day-of-week median at every intermittent node where both scores are
defined, so that rule would install a uniformly worse forecaster throughout. Adoption is therefore
decided on a validation block disjoint from the block the result is reported on, in the partition
of Figure~\ref{fig:blocks}. Disjointness is necessary and not sufficient, because a validation
score is itself an estimate with a standard error and a rule adopting on any margin adopts partly
on noise. Adoption therefore carries a one-standard-error margin, the device introduced for tree
pruning by \citet{breiman_classification_1984}: the intermittent estimator displaces the
day-of-week median only where its mean advantage exceeds one standard error of that advantage,
the paired differential taken over disjoint seven-day sub-blocks of the validation window.
Sub-blocks rather than days, because daily errors here are serially correlated and a daily
standard error would be biased downward; disjoint rather than overlapping, so no bootstrap is
needed to undo that correlation; and seven days because that is the dispersion unit already used
elsewhere. The rule is pre-registered and fails closed in favour of the incumbent on fewer than
two sub-blocks, on a non-finite differential and on zero dispersion. What it refuses is an
advantage small relative to its own variability, not one merely small. Appendix~D reports the
classification's sensitivity to the constant pair.

\section{Exogenous covariates and availability lead}
\label{sec:exo}

Two rungs of the ladder and one arm of the interval study take covariates beyond the target
series' own history. The set has seven members in three groups: weather supplies temperature,
rainfall, sunshine hours and a dry-day indicator thresholded at one millimetre; institutional
calendar supplies school-term and university-term indicators, which matter in a university city
where much of both trade and staffing follows the academic year; and events supply a
nearby-fixture indicator scoped per venue, so an anchor in one town cannot reach a venue in
another.

Membership is governed by what is knowable at the origin rather than by what the store holds, and
that is the substance of this section. A covariate is usable at serving only if a value for the
target date is obtainable when the forecast is issued. Calendar and event members satisfy this by
construction, since term dates and fixture lists are published in advance. Weather does not: a
weather value for a date seven days ahead is a forecast and carries the error of one, whereas the
value most readily available in a research store is the one recorded after the fact. Scoring an
exogenous model on recorded weather and serving it on forecast weather measures a system that
cannot be deployed. One feature is excluded on the same principle: a discount-share indicator
marking days on which nearly all takings were discounted is computed and used retrospectively to
flag anomalous days, and is never offered as a forward regressor, since its value on a future day
is not knowable at the origin. Admitting it would make the exogenous path a leak.

The basis is therefore fixed rather than assumed, and five arms of one exogenous entrant differ in
nothing but the weather handed to them for the seven-day window. Arm~N receives none and
reproduces the univariate rung, making it the reference the others are read against; arm~O
receives reanalysis at the valid time, which is perfect foreknowledge and so an upper bound rather
than a candidate; arm~H receives the archived value near the valid time; arm~F the forecast issued
at a fixed three-day operational lead; and arm~M the forecast issued $h$ days ahead for horizon
step $h$, which is serving reality and the only arm a deployed system could reproduce. Arms~F
and~M draw from one pinned global weather model across all leads, so their difference isolates the
lead policy rather than confounding it with a change of vendor model, and the maximum lead is
capped at the seven-day horizon so no arm receives a covariate for a date beyond the range the
forecast is issued for. Results are reported per horizon step as well as pooled, because forecast
skill decays with lead and a pooled average hides the gap where it is largest.

\section{Candidate models and the adoption gate}
\label{sec:ladder}

Model choice proceeds up a ladder of five rungs, each admitting structure the rung below cannot
represent, so capacity is bought one property at a time. The five supply nine scored entrants.
Rung 0 is a seasonal-naive benchmark.
Rung 1 is a robust day-of-week baseline taking a per-weekday median scaled by a monthly index
clipped to $[0.5, 2.0]$ and a bank-holiday factor, and fits nothing. Rung 2 covers exponential
smoothing with additive trend and additive seasonality at period 7, and a robust seasonal-trend
decomposition at the same period. Rung 3 is histogram gradient boosting forecasting recursively
on its own lag features, in a per-venue and a global arm. Rung 4 is a pretrained time-series
foundation model \citep{ansari_chronos_2024}, zero-shot and without per-venue training, in a
univariate arm and in an arm conditioned on the exogenous set of Section~\ref{sec:exo}; its
served point forecast is the median quantile clipped at zero.

A rung is adopted only if it beats both the seasonal-naive benchmark and the robust day-of-week
baseline on rolling-origin root mean squared scaled error at the basis ruled in
Section~\ref{sec:ruler}, and at Ellel on the unscaled equivalent. The ladder exists to make the
null expensive to reject: a foundation model is served only where it defeats a benchmark that
costs nothing to compute. Appendix~C states the adoption rule and its fail-closed conditions as
pseudocode.

% Trace: brain/log/45_G17d_Intermittency_and_Occurrence.md (frame lengths at ceiling 2026-07-07).
The exponential-smoothing rung is one fixed specification, not a member selected from the family
by an automatic procedure. The information-criterion search over trend, seasonal and damping
options that standard implementations run by default is deliberately not performed: at frame
lengths of 399 days at the longest it would select among specifications on a sample that cannot
separate them, and would add a second selection step whose uncertainty the rolling-origin
comparison does not carry. A poor showing by that rung is therefore evidence against the
specification and not against exponential smoothing as a class.

% Trace: brain/log/68_R5_tabpfn_entrant_result.md (pre-registered abort conditions).
Two further entrants are specified and scored at no venue, for reasons outside the data. One has
no backend in the evaluation environment.
The other, a tabular-backbone forecaster, routes inference through a vendor service in its
released configuration, and this estate's revenue may not leave the machine; a local route raises
a separate licensing question that Section~\ref{sec:further-work} carries. The estate is refitted
weekly, with an additional refit on a confirmed change point, and that policy is what makes the
protocol of Section~\ref{sec:injection} necessary.

\section{Model comparison procedure}
\label{sec:selection}

Comparing rung means over six rolling origins covering 42 consecutive days and adopting the
argument-minimum is not a procedure this evidence supports, for a reason that is arithmetic
rather than a matter of degree. The \citet{diebold_comparing_1995} test of equal predictive
accuracy is oversized in small samples, and the correction of \citet{harvey_testing_1997}
multiplies the statistic by
\begin{equation}
\left[\frac{n + 1 - 2h + n^{-1}h(h-1)}{n}\right]^{1/2},
\label{eq:hln}
\end{equation}
for $n$ forecast comparisons at horizon $h$. At $n = 6$ and $h = 7$ the numerator is
$6 + 1 - 14 + 7 = 0$, so the factor is exactly zero and no corrected statistic is computable on
any data whatever. Such a design does not have a weak test; it has none.
% Trace: brain/log/63_R4_metric_ordering_result.md.

Advancing the origin by one day rather than by the horizon lifts the Beer Hall from 39 origins to
273 and raises Equation~\ref{eq:hln} to 0.976, giving 273, 260 and 205 origins at the Beer Hall,
Ellel and Two River Taps. Overlapping test windows are the point of the change and carry a cost
handled rather than ignored: consecutive seven-day windows at a one-day step share six of seven
observations, so loss differentials are serially correlated by construction and are not
independent draws. Figure~\ref{fig:origins} in Appendix~C draws the scheme with a magnified inset
in which that overlap is visible.

\label{sec:mcs}
With a computable test available the selection question is restated. Rather than which rung has
the lowest mean, this work asks which rungs the data cannot separate, using the model confidence
set of \citet{hansen_model_2011}, in which the hypothesis of equal predictive ability is tested
repeatedly across the candidate set, one model eliminated at a time, until it is no longer
rejected. All nine entrants enter at every venue with no venue-specific shortlist; the statistic
is the range statistic with its paired elimination rule; resampling is a moving block bootstrap at
block length seven, chosen to respect the serial correlation the one-day step induces, with
$B = 1000$ replications, seed 93 and $\alpha = 0.10$ primary. It runs on the folds common to all
nine entrants, which need not be the full origin counts above and are recorded per venue in
Table~\ref{tab:mcs-config} in Appendix~C, with the full configuration, written to the project's
decision log before the procedure ran.

Two properties govern how the output is read. The procedure is conservative by construction,
controlling error across the whole candidate set, so a difference a pairwise test detects can
still sit inside the set, and where that occurs both facts are reported. And a large retained set
is a valid outcome rather than a failure: \citet{hansen_model_2011} are explicit that where the
data lack the information to separate models the procedure terminates early and returns many,
which is a statement about the evidence rather than a defect of the test. Because the compared
models are strongly correlated across folds, the variance of a paired loss differential can be far
smaller than the marginal standard errors suggest; that quantity is measured rather than assumed,
and Appendix~D reports it with a block-length sweep.

\section{Conformal interval construction}
\label{sec:conformal}

A deviation detector needs a normal range, which a point forecast does not supply. Intervals are
constructed by split conformal prediction: distribution-free, and giving finite-sample coverage
without a Gaussian assumption, which daily revenue does not admit since it is skewed and bounded
below. Given absolute residuals on a calibration set of size $n$, the half-width at level
$1-\alpha$ is the $k$-th smallest absolute residual with $k = \lceil (n+1)(1-\alpha) \rceil$, and
the interval is the point forecast plus and minus that quantity. Calibration is refreshed on
rolling seven-day blocks from residuals strictly earlier than the block, in the temporally robust
spirit of \citet{xu_conformal_2021}; Appendix~C states the construction and its group-conditional
variant as pseudocode.

% Figure fig:blocks (out/fig_blocks.tex) is inserted here.
Figure~\ref{fig:blocks} shows how the calendar is partitioned to support it, and the partition
buys the guarantee itself rather than tidiness. A forecaster scored against the points it was
fitted to produces an optimistic residual, and a quantile of that residual carries no coverage
claim. Every conformal band reported here is built on the lower arrangement in that figure, the
same fitted forecaster producing both the calibration scores and the forecasts the band is
attached to.

That index $k$ can exceed $n$. Where $n < (1-\alpha)/\alpha$, which is fewer than nine calibration
points at 90 per cent, the correct conformal set is unbounded, so the guarantee is unavailable at
that
sample size rather than merely weak. Reference implementations return an infinite interval; this
work returns the largest observed residual, because an unbounded interval is useless to an
operator deciding whether a day is normal. The substitution is handled by counting rather than by
silence: every band issued below the attainable size is tallied beside the coverage it qualifies,
since a guarantee lapsing quietly would lapse first on the thinnest groups.

Coverage is bounded above as well as below. \citet{angelopoulos_conformal_2023} give expected
coverage as at most $1-\alpha + 1/(n+1)$, so a coverage materially above nominal signals
misspecification as much as one below, and both directions are tested per venue. That bound also
assumes almost surely distinct conformity scores, which a closure calendar violates: doors shut
gives an actual and a forecast of zero and a residual of exactly zero, an atom at the foot of the
score distribution. The condition is tested rather than assumed.

Three refinements are specified. A Mondrian variant computes group-conditional quantiles
separating active from structural-zero days, so a closed venue's near-zero residuals cannot shrink
a trading day's interval. Per-step calibration fits a separate quantile per horizon step, since a
one-step and a seven-step residual are not exchangeable. And two adaptive methods are implemented:
adaptive conformal inference \citep{gibbs_adaptive_2021}, which updates the effective miscoverage
level online after each outcome, and its aggregated form \citep{zaffran_adaptive_2022}, which runs
one learner per rate on a grid and aggregates online, which removes the learning-rate choice the
former exposes. Aggregation is applied separately to each bound under its own pinball loss, so a
miss on one cannot reweight the other.

The Mondrian partition is an observed calendar variable rather than an inferred regime, and its
three claims come from three places of which only one is a guarantee. What the construction
\emph{is} is described by \citet{barber_conformal_2023}: groups are formed and observations
assumed exchangeable within each. What it \emph{guarantees} is the split-conformal property
applied within each group, which \citet{stocker_gentle_2025} state as finite-sample marginal
coverage at the nominal level. Why the partition is observed rather than inferred is where care is
owed. \citet{sun_conformal_2025} bound a state-conditional method's coverage error by a term
scaling with the rate at which its state predictions are wrong, so a method obliged to infer its
regime pays a penalty a known calendar does not. That motivates the choice and does not certify
it: the result is proved for an algorithm carrying an adaptive miscoverage update this work does
not implement, and citing it as the guarantee would claim a theorem about a procedure the served
system never runs.

The adaptive alternative was not left as an argument. Both methods were implemented and measured
against the static variants at this estate's one genuine regime change, and the band that is
served is the one that outscored the others there; Section~\ref{sec:res-winkler} reports the
comparison. Direct evidence from the series in question is a stronger warrant for a design than a
citation, which is why the choice was settled by measurement. Interval methods are compared on the
Winkler score as primary, with
coverage and mean width beside it, because comparing on coverage alone rewards unlimited widening
while the score penalises width and miscoverage together.

\section{Deviation detection}
\label{sec:detection}

Detection is separated from explanation by design. The statistical layer reports what it observed
and returns null where it has no feed for a cause; it never attributes. A separate enrichment
layer, on meeting a flagged null, retrieves candidate coincident factors and reports them as
coincidences with evidence attached, and never rewrites detection. That separation is what allows
a flagged day with no available explanation to be reported honestly as unexplained.

A point deviation is a band breach. The standardised residual is
\begin{equation}
z_t \;=\; \frac{y_t - \tilde{y}_t}{\max\bigl(w_t^{(90)},\, \varepsilon\bigr)},
\label{eq:z}
\end{equation}
with $\tilde{y}_t$ the day-of-week median and $w_t^{(90)}$ the conformal half-width at 90 per
cent, so $\lvert z_t \rvert = 1$ sits exactly on the interval edge and a deviation is by
construction a breach. Tying the threshold to the interval means detection inherits the coverage
properties of Section~\ref{sec:conformal}. A sustained shift needs a different instrument, since
one large day should not raise a regime alarm. A cumulative sum \citep{page_continuous_1954} of
the form $S^{+}_t = \max(0,\, S^{+}_{t-1} + z_t - k)$ with $k = 0.5$ alarms at $S^{+} > 5$,
two-sided: a day at $z = 3$ contributes only $2.5$, whereas ten consecutive days near $z = 1$
reach the threshold. A persistence rule alarms on four same-direction breaches within seven
trailing trading days, catching an abrupt level shift the cumulative sum reaches more slowly.
Bayesian online change-point detection \citep{adams_bayesian_2007} is implemented as a benchmark
and does not drive production alarms. Preferring an online cumulative sum to offline segmentation
follows the distinction drawn by \citet{truong_selective_2020}: only the online family is
admissible where a decision must be taken on the day. Appendix~C states both detectors as
pseudocode.

What that pairing does and does not inherit from its literature is stated plainly, because the
constants are borrowed and the unit is not. The operating point $k = 0.5$, $h = 5$ and the
average-run-length tables recommending it are derived for a statistic accumulating deviations
standardised by a \emph{standard deviation}. Equation~\ref{eq:z} does not divide by a standard
deviation. It divides by a conformal half-width, a quantile of the absolute residual, and the
ratio between the two depends on the shape of the residual distribution rather than being a fixed
constant. The constants are therefore transplanted into a unit in which their published
false-alarm properties are undefined, and no average-run-length claim is made for this detector on
their authority. What licenses the operating point instead is empirical validation against a real
closure and against synthetic injections at known onsets, on this estate's own residuals: a
weaker warrant than a closed-form run length, and the honest one. The alternative is to
restandardise $z_t$ by a robust scale estimate and recover the tables, at the cost of decoupling
the detector from the interval whose coverage Section~\ref{sec:conformal} establishes. The
coupling was judged worth more, since a deviation that is not a band breach would be a second and
separately calibrated notion of normal.

\section{Occurrence modelling}
\label{sec:occurrence}

A venue that is closed has no demand rather than zero demand, and a model treating the two alike
forecasts a fraction of a trading day's revenue onto a day the doors are shut. The standard
remedy is a two-part specification \citep{cragg_statistical_1971, mullahy_specification_1986}, in
which a binary model governs whether the outcome is positive and a second governs the amount:
\begin{equation}
\hat{y}_t \;=\; P(\text{trade at } t) \cdot \mathrm{E}\bigl[\,y_t \mid \text{trade at } t\,\bigr].
\label{eq:hurdle}
\end{equation}

% Trace: brain/log/67_DD3_hurdle_saturation_result.md (seed 93, n = 400).
The first factor is estimated rather than observed: it is the conditional mean of the occurrence
label given the day of the week over the training slice, a saturated nonparametric estimator with
one categorical covariate. It returns zero or one at the Beer Hall because that venue's closure
calendar is deterministic, not because the code imposes those values, and a venue trading on some
Mondays and not others would return a fraction. For a design matrix of day-of-week indicators the
maximum-likelihood estimate of a probit or logit first stage \emph{is} the within-cell empirical
frequency, and the two agree here to $7.6 \times 10^{-5}$, optimiser tolerance rather than a
modelling difference. What the closed form buys is numerical: in the deterministic cells the
fitted probability converges while the coefficient does not, a controlled check driving
$\lvert\hat{\beta}\rvert$ past 11 and still diverging at two thousand iterations, the signature of
complete separation.

What the design forgoes is covariates, not estimation. Both structural gains the specification is
motivated by survive: the amount model is fitted on trading days alone and so is not diluted by
structural-closure zeros, and the two processes are not constrained to share parameters. The
honest limitation is that this first stage conditions on the day of the week and nothing else,
adequate where closure is contractual and inadequate where it is event-led. That is Ellel, whose
booking diary was unavailable, so the occurrence covariate is left inert and the gate reported as
untestable there rather than approximated. The tempting approximation, deriving Ellel's occurrence
from its own trading days, is circular, and the implementation makes it impossible rather than
discouraged: the diary function takes no revenue parameter, and a test asserts no such argument or
branch exists.

\section{Detection evaluation protocol}
\label{sec:injection}

Detection performance is measured against a corpus of 644 synthetic events injected into the
history: sustained regime shifts, single-day spikes, stock drawdowns, and events coincident with
an exogenous factor. Recall and latency are scored against known onsets.
% Trace: brain/log/50_G17i_Injection_Realism.md.
That corpus perturbs the standardised residual stream: it displaces $z$, back-solves the actual
value while holding the forecast expectation fixed, and splices perturbed test rows onto
unperturbed training rows, so the forecaster never sees the perturbation in either its expectation
or its training data. Production refits weekly, refits again on a confirmed change point, and
recalibrates on rolling seven-day blocks, and all three absorb a sustained shift. Recall and
latency measured on that corpus are therefore upper bounds, and the size of the gap is an
empirical question rather than a matter for assertion.

That corpus is retained unmodified as the control arm, and a second pipeline supplies the
realistic one. The realistic arm fixes the perturbation in revenue units at the injection date,
applies it to the
raw revenue of every row from that date onward including rows that later become training data,
re-derives the forecast and the interval under the production refit policy, and only then runs the
unchanged detection code. Fixing the magnitude in revenue rather than in $z$ is what makes the
arms comparable: the same money moves in both, and the pipeline is the only difference. Injections
land only on trading days, using the occurrence definition of Section~\ref{sec:occurrence}. The
paired comparison runs on a stratified subsample of $n = 120$ under a recorded seed rather than on
all 644 events, because the realistic arm re-runs the forecaster for every event; Ellel is
excluded, since without the booking diary its occurrence label is inert and an injection there
could not be verified to land on a trading day. Appendix~D records the stratification and the
seed. The design also makes a feedback loop observable: a confirmed change point triggers a refit,
a refit deflates subsequent standardised residuals, and one detection can therefore suppress later
detection of the same continuing event.

\section{Knowledge-gap signal}
\label{sec:chatlog}

The chat corpus of Section~\ref{sec:design} is used to surface repeated operational questions as
candidate gaps in the estate's documented procedures. Messages are embedded and clustered, and a
cluster is surfaced only where its density of unanswered questions exceeds an empirically measured
baseline for ordinary traffic. The threshold protects the ranked output, since a gap detector
firing on routine chatter would train a reader to ignore the whole briefing; the signal is
therefore tuned for precision rather than recall. The embedding backend is pinned to a keyless
term-frequency projection rather than an alternative requiring a network service or a downloaded
model, because three backends produce three clusterings and a result depending on which was
configured would not be reproducible; the backend name is stamped on every artefact it writes.

This corpus differs from every other input in that its write path does not belong to the author.
Staff compose the messages the signal reads, so material reaching the briefing is supplied by the
people the briefing describes. \citet{zou_poisonedrag_2025} show what that admits: injecting a few
malicious texts into a retrieval-augmented system's knowledge database can induce a language model
to return an attacker-chosen answer. Three properties bound the exposure and none is a mitigation.
The corpus is single-tenant, so an attacker must already be staff; the forecast path does not read
it, so no injected passage can move a number; and the agent's verdicts are not served. What such a
passage can reach is the narrative layer of the briefing, which is the layer a reader has least
independent means to check. Section~\ref{sec:res-chatlog} reports what the signal returned on
this corpus.

\section{Intervention layer and scoring apparatus}
\label{sec:agent}

Whether a language-model agent can decide which detected deviations are worth raising is a
question the served artefact does not address, since it ranks candidates by a product of six
constants with no model participating in any served decision. This section specifies the
apparatus that addresses it. For each candidate item the agent returns a probability that the item
is worth raising, with a rationale and the hash of the prompt that produced it. Temperature is
zero and the model identifier pinned. Every response is cached on disk keyed on prompt, model and
item payload, and the evaluation replays from that cache without service calls, so a reader
without a credential can reproduce every number.

Two features of the protocol foreclose specific criticisms. The prompt is committed in a revision
provably preceding the one introducing any evaluation code, so that ordering establishes it was
fixed before any score existed. And the decision threshold is not elicited as a single number:
following \citet{fu_prism_2026}, the ratio of the cost of a missed problem to that of a false
alarm is swept over a pre-registered grid, so an operator's elicited ratio selects a point on a
curve already computed. Whether the agent contributes is decided against the incumbent
constant-weighted score and against a random baseline at the matched base rate, on the same corpus
and grid, with a paired bootstrap on decision quality; close agreement between agent and constants
is declared in advance to be a negative result rather than reframed.

Two limits bound what the apparatus can establish. The injection corpus supplies ground truth for
whether a deviation occurred, not for whether a manager should have been told, so what is measured
is \emph{detection} calibration and the code records that label. And no agreement statistic
against a model-based judge is reported: \citet{bavaresco_llms_2025} find such agreement varies
substantially across models and datasets and depends on whether the judged text is human- or
model-generated, concluding that a judge must be validated against human judgements before it is
relied upon. Absent that annotation a judge here would be an unvalidated proxy.
